% Faithful editable transcription: official article32 pp.32:2,32:5,32:8--32:12.
\section*{Author matrix conventions and complexity context}
One may view derandomizing BPL as the problem of approximating powers of substochastic matrices. For a TM with $s$ bits of memory, one can label all its states by elements in $[2^s]$. We can define $A\in\mathbb R^{2^s\times2^s}$ to be its transition matrix in the sense that $A_{i,j}$ is the probability that the machine moves from state $i$ to state $j$ in one step. Note that it must arrive at an accept or reject state in $2^s$ steps, so we only need to approximate $A^{2^s}$.
\paragraph{Definition 4 ($L_1$-norm).} Define the $L_1$-norm of a vector $(x_1,\ldots,x_n)^\top\in\mathbb R^n$ to be
\[\|(x_1,\ldots,x_n)^\top\|_1:=|x_1|+\cdots+|x_n|.\]
Define the $L_1$-norm of a matrix $A\in\mathbb R^{n\times n}$ to be
\[\|A\|_1:=\sup_{x\in\mathbb R^n}\frac{\|Ax\|_1}{\|x\|_1}=\max_{1\le j\le n}\{|A_{1,j}|+|A_{2,j}|+\cdots+|A_{n,j}|\}.\]
\paragraph{Theorem 5.} For any $A,B\in\mathbb R^{n\times n}$, we have:
\begin{enumerate}\item $\|A+B\|_1\le\|A\|_1+\|B\|_1$;\item $\|AB\|_1\le\|A\|_1\|B\|_1$;\item If $\|A\|_1,\|B\|_1\le1$, then for any $p\in\mathbb Z^+$, $\|A^p-B^p\|_1\le p\|A-B\|_1$.\end{enumerate}
\paragraph{Definition 6 (Non-negative Matrix).} We say a matrix is non-negative if each of its entry is non-negative.
\paragraph{Definition 7 (Substochastic Matrix).} We say a matrix $A\in\mathbb R^{n\times n}$ is a substochastic matrix if $A$ is non-negative and $\|A\|_1\le1$.
For simplicity, we always assume that the size of a substochastic matrix is a power of 2. To represent a substochastic matrix, we independently represent each entry in binary, accurate to $100\log n$ decimal places.
% STATEMENT_BEGIN
\paragraph{Theorem 20.} Let $n$ be a power of 2. Then there exists a $\operatorname{poly}(n)$-size $O(\log n)$-depth L-uniform AC circuit family $\{\mathcal M_n\}$\footnote{We require that given $n$, description of $\mathcal M_n$ can be computed in space $O(\log n)$.} such that on input a substochastic matrix $A\in\mathbb R^{n\times n}$, $\mathcal M_n$ outputs a substochastic matrix $M\in\mathbb R^{n\times n}$ such that $\|M-A^n\|_1\le1/n$.
\paragraph{Corollary 21.} $\mathrm{BPL}\subseteq\mathrm{L}\text{-}\mathrm{AC}^1$.
% STATEMENT_END
% PROOF_BEGIN
\section*{4 The Iteration Method}
In this section, we will introduce the iteration step, which is the core of our proof of $\mathrm{BPL}\subseteq\mathrm{L}\text{-}\mathrm{AC}^1$.
\paragraph{Theorem 18 (The Iteration).} Let $A\in\mathbb R^{n\times n}$ be a substochastic matrix and $k,t\in\mathbb Z^+$ such that $\log n\ge k\ge t$. Suppose substochastic matrices $B_0,\ldots,B_{k-1}$ are approximations of $A^{2^0},\ldots,A^{2^{k-1}}$ such that $\|B_i-A^{2^i}\|_1\le\varepsilon_i$ for $i=1,2,\ldots,k-1$. Define
\begin{align*}
C:=&-\sum_{i=1}^{t-1}\sum_{\substack{\{j_1<\cdots<j_p\}\uplus\{j_1'<\cdots<j_q'\}\\=\{k-1,k-2,\ldots,k-i+1\}}}B_{j_p}\cdots B_{j_1}B_{k-i}^2 B_{j_1'}\cdots B_{j_q'}\\
&+\sum_{\substack{\{j_1<\cdots<j_p\}\uplus\{j_1'<\cdots<j_q'\}\\=\{k-1,k-2,\ldots,k-t+1\}}}B_{j_p}\cdots B_{j_1}B_{k-t}^2 B_{j_1'}\cdots B_{j_q'}.
\end{align*}
Then
\[\|C-A^{2^k}\|_1\le\sum_{i=1}^{t-1}2^{i-1}\varepsilon_{k-i}^2+2^t\varepsilon_{k-t}.\]
\begin{proof}
Note that
\begin{align*}
C-A^{2^k}=&-\sum_{i=1}^{t-1}\sum_{\substack{\{j_1<\cdots<j_p\}\uplus\{j_1'<\cdots<j_q'\}\\=\{k-1,k-2,\ldots,k-i+1\}}}B_{j_p}\cdots B_{j_1}(A^{2^{k-i}}-B_{k-i})^2B_{j_1'}\cdots B_{j_q'}\\
&-\sum_{\substack{\{j_1<\cdots<j_p\}\uplus\{j_1'<\cdots<j_q'\}\\=\{k-1,k-2,\ldots,k-t+1\}}}B_{j_p}\cdots B_{j_1}(A^{2^{k-t+1}}-B_{k-t}^2)B_{j_1'}\cdots B_{j_q'}.
\end{align*}
So by Theorem 5,
\begin{align*}\|C-A^{2^k}\|_1&\le\sum_{i=1}^{t-1}2^{i-1}\|A^{2^{k-i}}-B_{k-i}\|_1^2+2^t\|A^{2^{k-t}}-B_{k-t}\|_1\\&\le\sum_{i=1}^{t-1}2^{i-1}\varepsilon_{k-i}^2+2^t\varepsilon_{k-t}.\end{align*}
\end{proof}
\paragraph{Theorem 19 (Computing the Iteration).} Let $n,k,t,A,B_0,\ldots,B_{k-1},\varepsilon_0,\ldots,\varepsilon_{k-1},C$ be as defined in Theorem 18. Let $d$ be an integer such that $4\log n\ge d\ge t/10$. Then there exists a $\operatorname{poly}(n)$-size $O(d)$-depth L-uniform AC circuit family $\{\mathcal I_{n,k,t,d}\}$ such that on inputs $B_{k-t},\ldots,B_{k-1}$, if
\[\sum_{i=1}^{t-1}2^{i-1}\varepsilon_{k-i}^2+2^t\varepsilon_{k-t}\le\frac1{2^{d+2}}\]
is satisfied, then $\mathcal I_{n,k,t,d}$ outputs a substochastic matrix $C'$ such that $\|C'-A^{2^k}\|_1\le1/2^d$.
The intuition behind Theorem 19 is that to approximately compute $C$, all arithmetic operations only need a multiplicative accuracy of $1/2^{\Theta(d)}$. This can be done efficiently by L-uniform AC circuit by Theorem 14.
\begin{proof}[Proof of Theorem 19]
We observe that $C$ is the sum of $2^{t-1}$ ``+'' terms and $2^{t-1}-1$ ``-'' terms, and each term is a multiplication of not more than $t+1$ substochastic matrices. We will first show how to approximate the multiplication of substochastic matrices and then show how to approximate their sum.
To approximate $Z:=XY$ for two substochastic matrices $X,Y$, we only need to approximate $\sum_{r=1}^n X_{i,r}Y_{r,j}$ for each pair $(i,j)\in[n]^2$. We first represent each entry $X_{i,r},Y_{r,j}$ using $n^{100}$ bits such that fraction of 1's in these $n^{100}$ bits is equal to the entry, then use a layer of AND gate to represent each $X_{i,r}Y_{r,j}$ using fraction of 1's in $n^{200}$ bits, and then represent each $\frac1n\sum_{r=1}^n X_{i,r}Y_{r,j}$ using fraction of 1's in $n^{201}$ bits. Then we invoke $C_{n^{201},\ell,\lceil\ell(1+1/2^{20d+10})\rceil}$ (as defined in Theorem 14, which has depth $\le O(\frac d{\log\log n}+1)\le O(\frac d{\log(t+1)})$)\footnote{In Theorem 14 we take $(a,b)=(\ell,\lceil\ell(1+1/2^{20d+10})\rceil)$, and then $\log\frac b{b-a}\le O(d)$.} for $\ell=1,2,\ldots,n^{200}$ over these $n^{201}$ bits. Suppose $\ell_0$ is the smallest index such that $C_{n^{201},\ell_0,\lceil\ell_0(1+1/2^{20d+10})\rceil}$ outputs 0, then we have
\[\frac{\ell_0-1}{n^{200}}<Z_{i,j}<\frac{\ell_0(1+\frac1{2^{20d+10}})}{n^{200}}\]
and thus\footnote{Since $n^{200}Z_{i,j}$ is an integer, we have $\frac{\ell_0-1}{n^{200}}<Z_{i,j}\Longrightarrow\frac{\ell_0}{n^{200}}\le Z_{i,j}$.}
\[\frac{Z_{i,j}}{1+\frac1{2^{20d+10}}}-\frac1{n^{100}}\le\frac1{n^{100}}\left[\frac{\ell_0}{n^{100}}\right]\le Z_{i,j}.\]
Use $[\ell_0/n^{100}]/n^{100}$ as an approximation of $Z_{i,j}$, then we obtain an approximation $\widetilde Z$ of $Z$ such that $Z-\widetilde Z$ is non-negative and $\widetilde Z$ is substochastic and $\|Z-\widetilde Z\|_1\le1/2^{20d+10}+1/n^{99}$. We need to be careful that here we need a multiplicative small error on each entry and thus we need to strengthen Lemma 17 to Theorem 14.
Then multiplication of not more than $t+1$ substochastic matrices can be computed via $O(\log(t+1))$ layers of multiplication of two matrices. Recall that multiplying two matrices uses $O(\frac d{\log(t+1)})$-depth and has additive error $1/2^{20d+10}+1/n^{99}$. So the total depth for computing multiplication of not more than $t+1$ substochastic matrices is $O(d)$ and the total error is $\le t(1/2^{20d+10}+1/n^{99})\le1/2^{19d+5}$.
To summarize, suppose $C=-\sum_{i=1}^{2^{t-1}-1}D_i+\sum_{i=1}^{2^{t-1}}D_i'$, here each $D_i,D_i'$ is multiplication of some substochastic matrices. Then we can compute their approximations $\widetilde D_i,\widetilde D_i'$ in $O(d)$ depth such that $\|D_i-\widetilde D_i\|_1\le1/2^{19d+5}$ and $\|D_i'-\widetilde D_i'\|_1\le1/2^{19d+5}$.
We approximate $\frac1{2^{t-1}}\sum_{i=1}^{2^{t-1}-1}\widetilde D_i$ and $\frac1{2^{t-1}}\sum_{i=1}^{2^{t-1}}\widetilde D_i'$. Use the similar idea as summing $\frac1n\sum_{r=1}^n X_{i,r}Y_{r,j}$, we can compute substochastic matrices $C^-,C^+$ using $O(d)$-depth, such that
\begin{align*}\left\|C^--\frac1{2^{t-1}}\sum_{i=1}^{2^{t-1}-1}\widetilde D_i\right\|_1&\le\frac1{2^{19d+5}},\\\left\|C^+-\frac1{2^{t-1}}\sum_{i=1}^{2^{t-1}}\widetilde D_i'\right\|_1&\le\frac1{2^{19d+5}}.\end{align*}
Then $2^{t-1}(C^+-C^-)$ is a good approximation of $A^{2^k}$ since
\begin{align*}\|2^{t-1}(C^+-C^-)-A^{2^k}\|_1&\le2^{t-1}\left\|C^--\frac1{2^{t-1}}\sum_{i=1}^{2^{t-1}-1}\widetilde D_i\right\|_1+2^{t-1}\left\|C^+-\frac1{2^{t-1}}\sum_{i=1}^{2^{t-1}}\widetilde D_i'\right\|_1\\&\quad+\sum_{i=1}^{2^{t-1}-1}\|D_i-\widetilde D_i\|_1+\sum_{i=1}^{2^{t-1}}\|D_i'-\widetilde D_i'\|_1\\&\quad+\left\|-\sum_{i=1}^{2^{t-1}-1}D_i+\sum_{i=1}^{2^{t-1}}D_i'-A^{2^k}\right\|_1\\&\le\frac{2^{t-1}}{2^{19d+5}}+\frac{2^{t-1}}{2^{19d+5}}+\frac{2^{t-1}}{2^{19d+5}}+\frac{2^{t-1}}{2^{19d+5}}+\|C-A^{2^k}\|_1\\&\le\frac1{2^{9d+4}}+\left(\sum_{i=1}^{t-1}2^{i-1}\varepsilon_{k-i}^2+2^t\varepsilon_{k-t}\right)\\&\le\frac1{2^{9d+4}}+\frac1{2^{d+2}}.\end{align*}
Here the last step is from the statement of Theorem 19.
Finally we compute a substochastic matrix $C'$ which is a good approximation of $A^{2^k}$ and $2^{t-1}(C^+-C^-)$. Here we need to be careful that $C$ and $2^{t-1}(C^+-C^-)$ are not necessarily non-negative or substochastic (but $A^{2^k}$ is guaranteed substochastic). Let
\begin{align*}C''_{i,j}&:=\max\{2^{t-1}(C_{i,j}^+-C_{i,j}^-),0\},\\C'_{i,j}&:=\frac1{n^{100}}\left[C''_{i,j}\left(1-\frac1{2^{d+1}}\right)\cdot n^{100}\right].\end{align*}
We can compute $C'$ given $C^+,C^-$ by hardwiring the map $(C_{i,j}^+,C_{i,j}^-)\mapsto C_{i,j}'$, which is L-uniform. Obviously $C'$ is non-negative. Note that $C''$ is entrywise closer to $A^{2^k}$ than $2^{t-1}(C^+-C^-)$ and hence
\[\|C''-A^{2^k}\|_1\le\|2^{t-1}(C^+-C^-)-A^{2^k}\|_1\le\frac1{2^{9d+4}}+\frac1{2^{d+2}}.\]
Therefore $C'$ is substochastic since
\[\|C'\|_1\le\left(1-\frac1{2^{d+1}}\right)\|C''\|_1\le\left(1-\frac1{2^{d+1}}\right)\left(1+\frac1{2^{9d+4}}+\frac1{2^{d+2}}\right)\le1.\]
Also note that
\begin{align*}\|C'-A^{2^k}\|_1&\le\|C'-C''\|_1+\|C''-A^{2^k}\|_1\\&\le\frac1{n^{99}}+\frac1{2^{d+1}}\|C''\|_1+\|C''-A^{2^k}\|_1\\&\le\frac1{n^{99}}+\frac1{2^{d+1}}\left(1+\frac1{2^{9d+4}}+\frac1{2^{d+2}}\right)+\frac1{2^{9d+4}}+\frac1{2^{d+2}}\\&\le\frac1{2^d}.\end{align*}
To summarize, we can output a valid $C'$ in $O(d)$-depth. And the circuit is $\operatorname{poly}(n)$-size and L-uniform.
\end{proof}
\section*{5 The Complete Algorithm}
\begin{proof}[Proof of Theorem 20]
Only consider $\log n\ge10$. For $k,t\in\mathbb N$ such that $k\le\log n$ and $1\le t\le3\log n-2k$, we wish to compute a substochastic matrix $M(k,t)$, which is an approximation of $A^{2^k}$, such that $\|M(k,t)-A^{2^k}\|_1\le1/2^t$. Then $M:=M(\log n,\log n)$ is the desired matrix.
For $k=0$, we can trivially let $M(0,t):=A$. Now we show how to recursively compute $M(k_0,t_0)$ for $k_0=1,2,\ldots,\log n$.
In Theorem 18, take the same $n,A$ and take $k:=k_0$, take $B_{k-i}:=M(k-i,[t_0/2]+2i)$ for $1\le i\le k$. Then we can take $\varepsilon_{k-i}:=1/2^{[t_0/2]+2i}$ for $1\le i\le k-1$ and $\varepsilon_0=0$. Now we will invoke Theorem 18,19 by choosing parameter $t$ properly according to the following two cases.
\paragraph{Case 1. $k\le2t_0+2$.}
Take the parameter $t$ in Theorem 18 to be $t:=k$. Then
\[\sum_{i=1}^{k-1}2^{i-1}\varepsilon_{k-i}^2+2^k\varepsilon_0=\sum_{i=1}^{k-1}\frac1{2^{2[t_0/2]+3i+1}}\le\frac1{2^{t_0+2}}.\]
In Theorem 19 take $d:=t_0$. It is easy to verify that $\log n\ge k\ge t$ and $4\log n\ge d\ge t/10$ hold when we invoke Theorem 18,19. Given $B_{k-1},\ldots,B_0$, use $\mathcal I_{n,k_0,k_0,t_0}$ (defined in Theorem 19) we can compute a substochastic matrix $C'$ such that $\|C'-A^{2^k}\|_1\le1/2^{t_0}$.
\paragraph{Case 2. $k\ge2t_0+3$.}
Take $t:=2t_0+3$ in Theorem 18. Then
\[\sum_{i=1}^{2t_0+2}2^{i-1}\varepsilon_{k-i}^2+2^{2t_0+3}\varepsilon_{k-2t_0-3}\le\sum_{i=1}^{2t_0+2}\frac1{2^{2[t_0/2]+3i+1}}+\frac1{2^{[t_0/2]+2t_0+3}}\le\frac1{2^{t_0+2}}.\]
In Theorem 19 take $d:=t_0$. Given $B_{k-1},\ldots,B_0$, use $\mathcal I_{n,k_0,2t_0+3,t_0}$ we can compute a substochastic matrix $C'$ such that $\|C'-A^{2^k}\|_1\le1/2^{t_0}$.
To summarize, take $M(k_0,t_0):=C'$, we can compute $M(k_0,t_0)$ given $M(k_0-i,[t_0/2]+2i)$ for $1\le i\le k_0$, via a $\operatorname{poly}(n)$-size $O(t_0)$-depth L-uniform AC circuit.
Let $\gamma>0$ be a concrete constant such that we can compute $M(k_0,t_0)$ given $M(k_0-i,[t_0/2]+2i)$'s via a $\operatorname{poly}(n)$-size $\gamma t_0$-depth L-uniform AC circuit. Note that if $M(k_0-i,[t_0/2]+2i)$ can be computed in $2\gamma(2(k_0-i)+([t_0/2]+2i))$-depth for $1\le i\le k_0$, then $M(k_0,t_0)$ can be computed in
\[\gamma t_0+\max_{1\le i\le k_0}\{2\gamma(2(k_0-i)+([t_0/2]+2i))\}\le2\gamma(2k_0+t_0)\]
-depth. Also note that $M(0,t_0)$'s are just the inputs, so by induction we know $M(k_0,t_0)$ can be computed in $2\gamma(2k_0+t_0)$-depth. Specially, $M(\log n,\log n)$ (which is the desired output) can be computed in $6\gamma\log n\le O(\log n)$-depth. Also note that we use ``compute $M(k_0,t_0)$ given $M(k_0-i,[t_0/2]+2i)$'' $O((\log n)^2)$ many times, so the total circuit size for computing $M(\log n,\log n)$ is still $\operatorname{poly}(n)$.
\end{proof}
% PROOF_END
